% =========================================================================
% CAPÍTULO 10
% =========================================================================
\chapter{Comparativa Internacional y Trabajo Futuro}

La evidencia de esta tesis es consistente con la literatura internacional: los modelos de aprendizaje automático en panel no eliminan la utilidad de la econometría clásica, sino que la complementan, especialmente en presencia de heterogeneidad institucional y necesidad de generalización a entidades con historia corta.

\section{Trabajo Futuro}
\begin{enumerate}
    \item Sustituir las series macroeconómicas sintéticas por datos oficiales de la API de Banxico una vez disponible el token de acceso correspondiente, y reevaluar la importancia relativa de estas variables.
    \item Entrenar arquitecturas LSTM formales una vez que el panel acumule suficiente profundidad temporal por institución.
    \item Extender el horizonte de predicción más allá de un trimestre.
    \item Investigar específicamente el origen de los valores atípicos extremos en la variable NICAP directamente con la CNBV.
    \item Incorporar variables de gobernanza corporativa y análisis de sentimiento de prensa financiera.
\end{enumerate}

